\documentclass[tikz,border=6pt]{standalone}
\usepackage{ctex}
\usepackage{amsmath}
\usepackage{amssymb}
\usepackage{pgfplots}
\pgfplotsset{compat=1.18}
\usetikzlibrary{calc,angles,quotes,arrows.meta,positioning,decorations.markings}
\begin{document}
\begin{tikzpicture}
\begin{axis}[
  width=13cm,height=10cm,
  axis lines=middle,
  xlabel={$x$},ylabel={$y$},
  xlabel style={at={(axis description cs:1,0.5)},anchor=west},
  ylabel style={at={(axis description cs:0.5,1)},anchor=south},
  xmin=-2.6,xmax=2.6,
  ymin=-5.2,ymax=6.2,
  xtick={-2,-1,1,2},
  ytick={-4,-2,2,4,6},
  enlargelimits=false,
  clip=true,
  samples=200,
  domain=-2.5:2.5,
  legend pos=north west,
  legend cell align=left,
  legend style={font=\small,fill=white,fill opacity=0.85,text opacity=1,draw=black!40},
]
% e^x （细虚线，灰蓝）
\addplot[thick,gray!70!blue,dashed,domain=-2.5:2.5]{exp(x)};
\addlegendentry{$e^{x}$}
% e^{-x} （细虚线，灰红）
\addplot[thick,gray!70!red,dashed,domain=-2.5:2.5]{exp(-x)};
\addlegendentry{$e^{-x}$}
% cosh x = (e^x+e^-x)/2 （蓝，U 形，在两指数之上）
\addplot[very thick,blue,domain=-2.5:2.5]{(exp(x)+exp(-x))/2};
\addlegendentry{$\cosh x=\dfrac{e^x+e^{-x}}{2}$}
% sinh x = (e^x-e^-x)/2 （绿，奇函数）
\addplot[very thick,green!55!black,domain=-2.5:2.5]{(exp(x)-exp(-x))/2};
\addlegendentry{$\sinh x=\dfrac{e^x-e^{-x}}{2}$}
% 在 x=1.6 处画"半和/半差"竖直比较线
\def\xc{1.6}
% e^x 与 e^-x 在该处的点
\addplot[only marks,mark=*,mark size=1.6pt,gray!70!blue] coordinates {(\xc,{exp(\xc)})};
\addplot[only marks,mark=*,mark size=1.6pt,gray!70!red]  coordinates {(\xc,{exp(-\xc)})};
\addplot[only marks,mark=*,mark size=2.2pt,blue]  coordinates {(\xc,{(exp(\xc)+exp(-\xc))/2})};
\addplot[only marks,mark=*,mark size=2.2pt,green!55!black] coordinates {(\xc,{(exp(\xc)-exp(-\xc))/2})};
% 竖直细线串起三点，直观显示 cosh 在中点
\draw[black!55,thin] (axis cs:\xc,{exp(-\xc)}) -- (axis cs:\xc,{exp(\xc)});
\node[right,font=\footnotesize,blue] at (axis cs:\xc,{(exp(\xc)+exp(-\xc))/2}) {$\cosh$ 在中点};
% 半透明说明框，置于右下空白处，避开曲线峰值
\node[draw=black,fill=yellow!18,opacity=0.9,text opacity=1,
      rounded corners,align=left,font=\footnotesize]
  at (axis cs:-1.35,-3.4)
  {$\cosh x$：两指数的半和（偶，U 形）\\
   $\sinh x$：两指数的半差（奇）\\
   $e^x=\cosh x+\sinh x$};
\end{axis}
\end{tikzpicture}
\end{document}
